\noindent \textbf{16.02}

\section{Простейшие типы точек покоя}

\begin{itemize}
	\item Для анализа поведения решений линейных систем оказывается удобным выделить в картине интегральных кривых характерные элементы. Такими элементами оказываются окрестности точек покоя.
	
	\item Рассмотрим систему двух однородных ДУ с постоянными коэффициентами:
	\begin{equation} \label{eq:223}
		\begin{cases}
			\dv{x}{t} = a_{11}x + a_{12}y \\
			\dv{y}{t} = a_{21}x + a_{22}y
		\end{cases} \quad (223)
	\end{equation}
	причем $\Delta = \begin{vmatrix} a_{11} & a_{12} \\ a_{21} & a_{22} \end{vmatrix} \neq 0$
\end{itemize}

\begin{definition}{Точка покоя}{rest_point_223}
	Точка $x=x_0$, $y=y_0$, в которой правые части ур-ний системы \eqref{eq:223} обращаются в $0$, называется \textbf{точкой покоя} системы \eqref{eq:223}.
\end{definition}

Условие $\Delta \neq 0$ гарантирует, что точка покоя системы --- нулевая точка $x=0$, $y=0$.

\begin{itemize}
	\item Составим характеристическое ур-нение: 
	\begin{equation} \label{eq:224}
		\begin{vmatrix} a_{11} - \lambda & a_{12} \\ a_{21} & a_{22} - \lambda \end{vmatrix} = 0 \quad (224)
	\end{equation}
	и найдем его корни $\lambda_1, \lambda_2$. Возможны следующие случаи:
\end{itemize}

1) Корни $\lambda_1, \lambda_2$ характеристического ур-ния \eqref{eq:224} вещественные и разные:
\begin{itemize}
	\item[(а)] $\lambda_1 < 0, \lambda_2 < 0$. Точка покоя асимптотически устойчива (устойчивый узел) \textcolor{orange}{рис. 1}
	\item[(б)] $\lambda_1 > 0, \lambda_2 > 0$. Точка покоя неустойчива (неустойчивый узел). \textcolor{orange}{рис. 2}
	\item[(в)] $\lambda_1 > 0, \lambda_2 < 0$. Точка покоя неустойчива (седло). \textcolor{orange}{рис. 3}
\end{itemize}

\noindent\colorbox{green!30}{\parbox{\dimexpr\textwidth-2\fboxsep}{Траектории решений системы $(x(t), y(t))$ при изменении $t$ для различных типов точек покоя, на плоскости $xOy$:}}

\vspace{0.5cm}
\noindent
\begin{minipage}{0.48\textwidth}
	\centering
	\begin{tikzpicture}[scale=0.9, decoration={markings, mark=at position 0.6 with {\arrow{stealth}}}]
		\draw[-stealth] (-3,0) -- (3,0);
		\draw[-stealth] (0,-3) -- (0,3);
		\foreach \c in {0.5, 1.2, 2.0, 2.8} {
			\draw[postaction={decorate}] (2.5, \c) to[out=180, in=90] (0,0);
			\draw[postaction={decorate}] (-2.5, \c) to[out=0, in=90] (0,0);
			\draw[postaction={decorate}] (2.5, -\c) to[out=180, in=-90] (0,0);
			\draw[postaction={decorate}] (-2.5, -\c) to[out=0, in=-90] (0,0);
		}
		\draw[postaction={decorate}] (0,2.5) -- (0,0);
		\draw[postaction={decorate}] (0,-2.5) -- (0,0);
		\draw[postaction={decorate}] (2.5,0) -- (0,0);
		\draw[postaction={decorate}] (-2.5,0) -- (0,0);
	\end{tikzpicture}
	
	\footnotesize Рис. 1. Устойчивый узел. $\lambda_1, \lambda_2 \in \mathbb{R}$, $\lambda_1 \neq \lambda_2$, $\lambda_1 < 0, \lambda_2 < 0$. Асимптотически устойчива
\end{minipage}\hfill
\begin{minipage}{0.48\textwidth}
	\centering
	\begin{tikzpicture}[scale=0.9, decoration={markings, mark=at position 0.6 with {\arrow{stealth}}}]
		\draw[-stealth] (-3,0) -- (3,0);
		\draw[-stealth] (0,-3) -- (0,3);
		\foreach \c in {0.5, 1.2, 2.0, 2.8} {
			\draw[postaction={decorate}] (0,0) to[out=90, in=180] (2.5, \c);
			\draw[postaction={decorate}] (0,0) to[out=90, in=0] (-2.5, \c);
			\draw[postaction={decorate}] (0,0) to[out=-90, in=180] (2.5, -\c);
			\draw[postaction={decorate}] (0,0) to[out=-90, in=0] (-2.5, -\c);
		}
		\draw[postaction={decorate}] (0,0) -- (0,2.5);
		\draw[postaction={decorate}] (0,0) -- (0,-2.5);
		\draw[postaction={decorate}] (0,0) -- (2.5,0);
		\draw[postaction={decorate}] (0,0) -- (-2.5,0);
	\end{tikzpicture}
	
	\footnotesize Рис. 2. Неустойчивый узел. $\lambda_1, \lambda_2 \in \mathbb{R}$, $\lambda_1 \neq \lambda_2$, $\lambda_1 > 0, \lambda_2 > 0$
\end{minipage}

\vspace{0.5cm}
\begin{center}
	\begin{minipage}{0.6\textwidth}
		\centering
		\begin{tikzpicture}[scale=1, decoration={markings, mark=at position 0.5 with {\arrow{stealth}}}]
			\draw[-stealth] (-3,0) -- (3,0);
			\draw[-stealth] (0,-3) -- (0,3);
			\foreach \c in {0.3, 0.8, 1.5, 2.2} {
				\draw[postaction={decorate}] (\c, 2.8) .. controls (\c, \c) .. (2.8, \c);
				\draw[postaction={decorate}] (-\c, 2.8) .. controls (-\c, \c) .. (-2.8, \c);
				\draw[postaction={decorate}] (\c, -2.8) .. controls (\c, -\c) .. (2.8, -\c);
				\draw[postaction={decorate}] (-\c, -2.8) .. controls (-\c, -\c) .. (-2.8, -\c);
			}
			\draw[postaction={decorate}] (0,2.8) -- (0,0);
			\draw[postaction={decorate}] (0,-2.8) -- (0,0);
			\draw[postaction={decorate}] (0,0) -- (2.8,0);
			\draw[postaction={decorate}] (0,0) -- (-2.8,0);
		\end{tikzpicture}
		
		\footnotesize Рис. 3. Седло. $\lambda_1, \lambda_2 \in \mathbb{R}, \lambda_1 \neq \lambda_2, \lambda_1 > 0, \lambda_2 < 0$. Неустойчива
	\end{minipage}
\end{center}

\vspace{0.5cm}
3) Корни характеристического ур-ния \eqref{eq:224} комплексные: $\lambda_1 = p + iq$, $\lambda_2 = p - iq$
\begin{itemize}
	\item[(а)] $p < 0, q \neq 0$. Точка покоя асимптотически устойчива (устойчивый фокус) \textcolor{orange}{рис. 4}
	\item[(б)] $p > 0, q \neq 0$. Точка покоя неустойчива (неустойчивый фокус) \textcolor{orange}{рис. 5}
	\item[(в)] $p = 0, q \neq 0$. Точка покоя устойчива (центр) \textcolor{orange}{рис. 6}
\end{itemize}

\vspace{0.5cm}
\noindent
\begin{minipage}{0.48\textwidth}
	\centering
	\begin{tikzpicture}[scale=1, decoration={markings, mark=at position 0.4 with {\arrow{stealth}}, mark=at position 0.8 with {\arrow{stealth}}}]
		\draw[-stealth] (-2.5,0) -- (2.5,0);
		\draw[-stealth] (0,-2.5) -- (0,2.5);
		\draw[postaction={decorate}, domain=0:5*pi, samples=200, variable=\t] plot ({2.5*exp(-0.1*\t)*cos(-\t r)}, {2.5*exp(-0.1*\t)*sin(-\t r)});
	\end{tikzpicture}
	
	\footnotesize Рис. 4. Устойчивый фокус. $\lambda_1, \lambda_2 \in \mathbb{C}$, $\lambda_1 = \alpha + i\beta, \lambda_2 = \alpha - i\beta$, $\alpha < 0, \beta \neq 0$. Асимптотически устойчива
\end{minipage}\hfill
\begin{minipage}{0.48\textwidth}
	\centering
	\begin{tikzpicture}[scale=1, decoration={markings, mark=at position 0.4 with {\arrow{stealth}}, mark=at position 0.8 with {\arrow{stealth}}}]
		\draw[-stealth] (-2.5,0) -- (2.5,0);
		\draw[-stealth] (0,-2.5) -- (0,2.5);
		\draw[postaction={decorate}, domain=0:5*pi, samples=200, variable=\t] plot ({0.5*exp(0.1*\t)*cos(-\t r)}, {0.5*exp(0.1*\t)*sin(-\t r)});
	\end{tikzpicture}
	
	\footnotesize Рис. 5. Неустойчивый фокус. $\lambda_1, \lambda_2 \in \mathbb{C}, \lambda_1 = \alpha + i\beta, \lambda_2 = \alpha - i\beta$, $\alpha > 0, \beta \neq 0$.
\end{minipage}

\vspace{0.5cm}
\begin{center}
	\begin{minipage}{0.6\textwidth}
		\centering
		\begin{tikzpicture}[scale=1, decoration={markings, mark=at position 0.25 with {\arrow{stealth}}, mark=at position 0.75 with {\arrow{stealth}}}]
			\draw[-stealth] (-2.5,0) -- (2.5,0);
			\draw[-stealth] (0,-2.5) -- (0,2.5);
			\foreach \r in {0.5, 1.2, 2.0} {
				\draw[postaction={decorate}] (\r,0) arc (0:-360:\r);
			}
		\end{tikzpicture}
		
		\footnotesize Рис. 6. Центр. $\lambda_1, \lambda_2 \in \mathbb{C}, \lambda_1 = i\beta, \lambda_2 = -i\beta, \beta \neq 0$. Устойчива.
	\end{minipage}
\end{center}

\vspace{0.5cm}
3) Корни $\lambda_1 = \lambda_2$ кратные !
\begin{itemize}
	\item[(а)] $\lambda_1 = \lambda_2 < 0$. Точка покоя асимптотически устойчива (устойчивый узел) \textcolor{orange}{рис. 7, 8}
	\item[(б)] $\lambda_1 = \lambda_2 > 0$. Точка покоя неустойчива (неустойчивый узел) \textcolor{orange}{рис. 9, 10}
\end{itemize}

\vspace{0.5cm}
\noindent
\begin{minipage}{0.48\textwidth}
	\centering
	\begin{tikzpicture}[scale=0.9, decoration={markings, mark=at position 0.6 with {\arrow{stealth}}}]
		\draw[-stealth] (-2.5,0) -- (2.5,0);
		\draw[-stealth] (0,-2.5) -- (0,2.5);
		\foreach \c in {0.5, 1.2, 2.0} {
			\draw[postaction={decorate}] (2.5, \c) to[out=180, in=0] (0,0);
			\draw[postaction={decorate}] (2.5, -\c) to[out=180, in=0] (0,0);
			\draw[postaction={decorate}] (-2.5, \c) to[out=0, in=180] (0,0);
			\draw[postaction={decorate}] (-2.5, -\c) to[out=0, in=180] (0,0);
		}
		\draw[postaction={decorate}] (2.5,0) -- (0,0);
		\draw[postaction={decorate}] (-2.5,0) -- (0,0);
	\end{tikzpicture}
	
	\footnotesize Асимптотически устойчива. \\ Рис. 7. Точка покоя --- вырожденный устойчивый узел. $\lambda_1 = \lambda_2 = \lambda \in \mathbb{R}, \lambda < 0$.
\end{minipage}\hfill
\begin{minipage}{0.48\textwidth}
	\centering
	\begin{tikzpicture}[scale=0.9, decoration={markings, mark=at position 0.6 with {\arrow{stealth}}}]
		\draw[-stealth] (-2.5,0) -- (2.5,0);
		\draw[-stealth] (0,-2.5) -- (0,2.5);
		\foreach \angle in {30, 60, 120, 150, 210, 240, 300, 330} {
			\draw[postaction={decorate}] (\angle:2.5) -- (0,0);
		}
		\draw[postaction={decorate}] (0,2.5) -- (0,0);
		\draw[postaction={decorate}] (0,-2.5) -- (0,0);
		\draw[postaction={decorate}] (2.5,0) -- (0,0);
		\draw[postaction={decorate}] (-2.5,0) -- (0,0);
	\end{tikzpicture}
	
	\footnotesize Рис. 8. Устойчивый узел, вырожденный случай. $\lambda_1 = \lambda_2 = \lambda \in \mathbb{R}, \lambda < 0$. Асимптотически устойчива.
\end{minipage}

\vspace{0.5cm}
\noindent
\begin{minipage}{0.48\textwidth}
	\centering
	\begin{tikzpicture}[scale=0.9, decoration={markings, mark=at position 0.6 with {\arrow{stealth}}}]
		\draw[-stealth] (-2.5,0) -- (2.5,0);
		\draw[-stealth] (0,-2.5) -- (0,2.5);
		\foreach \c in {0.5, 1.2, 2.0} {
			\draw[postaction={decorate}] (0,0) to[out=0, in=180] (2.5, \c);
			\draw[postaction={decorate}] (0,0) to[out=0, in=180] (2.5, -\c);
			\draw[postaction={decorate}] (0,0) to[out=180, in=0] (-2.5, \c);
			\draw[postaction={decorate}] (0,0) to[out=180, in=0] (-2.5, -\c);
		}
		\draw[postaction={decorate}] (0,0) -- (2.5,0);
		\draw[postaction={decorate}] (0,0) -- (-2.5,0);
	\end{tikzpicture}
	
	\footnotesize Рис. 9. Неустойчивый узел. $\lambda_1 = \lambda_2 = \lambda \in \mathbb{R}, \lambda > 0$.
\end{minipage}\hfill
\begin{minipage}{0.48\textwidth}
	\centering
	\begin{tikzpicture}[scale=0.9, decoration={markings, mark=at position 0.6 with {\arrow{stealth}}}]
		\draw[-stealth] (-2.5,0) -- (2.5,0);
		\draw[-stealth] (0,-2.5) -- (0,2.5);
		\foreach \angle in {30, 60, 120, 150, 210, 240, 300, 330} {
			\draw[postaction={decorate}] (0,0) -- (\angle:2.5);
		}
		\draw[postaction={decorate}] (0,0) -- (0,2.5);
		\draw[postaction={decorate}] (0,0) -- (0,-2.5);
		\draw[postaction={decorate}] (0,0) -- (2.5,0);
		\draw[postaction={decorate}] (0,0) -- (-2.5,0);
	\end{tikzpicture}
	
	\footnotesize Рис. 10. Неустойчивый узел, вырожденный случай. $\lambda_1 = \lambda_2 = \lambda \in \mathbb{R}, \lambda > 0$.
\end{minipage}

\vspace{0.5cm}
\begin{itemize}
	\item \textbf{Пример 1}
\end{itemize}

\begin{example}{Пример 1}{ex1}
	Определим характер точки покоя $(0,0)$ системы:
	\begin{equation*}
		\begin{cases}
			\dv{x}{t} = 5x - y \\
			\dv{y}{t} = 2x + y
		\end{cases}
	\end{equation*}
	Составляем характеристическое ур-нение:
	\begin{equation*}
		\begin{vmatrix} 5-\lambda & -1 \\ 2 & 1-\lambda \end{vmatrix} = 0 \iff (5-\lambda)(1-\lambda) + 2 = 0
	\end{equation*}
	\begin{equation*}
		\lambda^2 - 6\lambda + 7 = 0
	\end{equation*}
	\begin{equation*}
		\begin{bmatrix}
			\lambda_1 = 3 + \sqrt{2} > 0 \\
			\lambda_2 = 3 - \sqrt{2} > 0
		\end{bmatrix}
	\end{equation*}
	корни вещественные, разные, положительные $\implies$ точка покоя $(0,0)$ --- неустойчивый узел.
\end{example}

\begin{itemize}
	\item Связь между типами точек покоя и коэффициентами характеристического ур-ния
\end{itemize}

Обозначения: $b = -(a_{11} + a_{22})$ \\
\hspace*{2.5cm} $\Delta = a_{11} \cdot a_{22} - a_{12} \cdot a_{21}$

Характеристич. ур-ние: 
\begin{align*}
	\begin{vmatrix} a_{11} - \lambda & a_{12} \\ a_{21} & a_{22} - \lambda \end{vmatrix} &= 0 \\
	\iff (a_{11} - \lambda)(a_{22} - \lambda) - a_{12}a_{21} &= 0 \\
	\iff a_{11}a_{22} - a_{11}\lambda - a_{22}\lambda + \lambda^2 - a_{12}a_{21} &= 0 \\
	\iff \lambda^2 - (a_{11} + a_{22})\lambda + a_{11}a_{22} - a_{12}a_{21} &= 0 \\
	\iff \text{\colorbox{yellow!30}{$\lambda^2 + b\lambda + \Delta = 0$}}
\end{align*}

Рассмотрим плоскость с прямоугольными декартовыми координатами $\Delta$ и $b$ и отметим на ней области, соответствующие различным типам точек покоя.

Из классификации следует, что условиями устойчивости точки покоя являются $\Re \lambda_1 < 0$ и $\Re \lambda_2 < 0$. Они выполняются при $\Delta > 0, b > 0$, для точек I четверти.

\vspace{0.5cm}
\noindent
\begin{minipage}{0.48\textwidth}
	\scriptsize
	\centering
	$\displaystyle[\begin{array}{l}
		\text{По т. Виета, } \Delta = \lambda_1 \lambda_2 > 0 \implies \text{корни одного знака} \\
		\lambda_1 + \lambda_2 = -b < 0 \implies \lambda_1 < 0, \lambda_2 < 0
	\end{array}]$
\end{minipage}\hfill
\begin{minipage}{0.48\textwidth}
	\scriptsize
	\centering
	$\displaystyle[\begin{array}{l}
		\lambda_1 = p+iq, \lambda_2 = p-iq, \quad p < 0 \\
		\Delta = \lambda_1 \lambda_2 = (p+iq)(p-iq) = p^2 + q^2 > 0 \\
		\lambda_1 + \lambda_2 = 2p = -b < 0 \implies b > 0
	\end{array}]$
\end{minipage}

\begin{center}
	$\downarrow$ \text{дискриминант } $\lambda^2 + b\lambda + \Delta = 0$
\end{center}

Если $\lambda_1, \lambda_2$ - комплексные, то есть $\mathcal{D} < 0$, точка покоя будет типа фокуса. Этому условию удовлетворяют точки, лежащие между ветвями $b^2 = 4\Delta$ и не принадлежат оси $O\Delta$ ($b^2 < 4\Delta, b \neq 0$).

\begin{center}
	\colorbox{cyan!30}{Типы точек покоя на $bO\Delta$}
	
	\vspace{0.5cm}
	\begin{tikzpicture}[scale=1.5]
		% Area Delta < 0
		\fill[orange!10] (-2, -3.2) rectangle (0, 3.2);
		
		% Area inside parabola
		\begin{scope}
			\clip (0,0) parabola (2.5, 3.16) -- (2.5, -3.16) parabola[bend at end] (0,0);
			\fill[cyan!10] (0,-3.5) rectangle (2.5, 3.5);
		\end{scope}
		
		% Area outside parabola (b > 0)
		\begin{scope}
			\clip (0,0) parabola (2.5, 3.16) -- (0, 3.16) -- cycle;
			\fill[blue!10] (0,0) rectangle (2.5, 3.5);
		\end{scope}
		
		% Area outside parabola (b < 0)
		\begin{scope}
			\clip (0,0) parabola (2.5, -3.16) -- (0, -3.16) -- cycle;
			\fill[blue!10] (0,0) rectangle (2.5, -3.5);
		\end{scope}
		
		% Axes
		\draw[-stealth, thick] (-2.5,0) -- (3,0) node[right] {$\Delta$};
		\draw[-stealth, thick] (0,-3.5) -- (0,3.5) node[above] {$b$};
		\node[below left] at (0,0) {$O$};
		
		% Parabola
		\draw[thick, orange, domain=0:2.5, samples=100] plot (\x, {2*sqrt(\x)}) node[right, black] {$b^2=4\Delta$};
		\draw[thick, orange, domain=0:2.5, samples=100] plot (\x, {-2*sqrt(\x)});
		
		% Labels
		\node[fill=orange!10, inner sep=1pt] at (-1.2, 1.5) {седла};
		
		\node[fill=cyan!10, inner sep=1pt] at (1.2, 1.2) {устойчивые фокусы};
		\node[fill=cyan!10, inner sep=1pt] at (1.2, -1.2) {неустойчивые фокусы};
		
		\node[rotate=45, fill=blue!10, inner sep=1pt] at (0.8, 2.5) {устойчивые узлы};
		\node[rotate=-45, fill=blue!10, inner sep=1pt] at (0.8, -2.5) {неустойчивые узлы};
	\end{tikzpicture}
\end{center}

Точки полуоси $b=0$, для которых $\Delta > 0$ соответствуют точкам покоя типа центра. Точки, расположенные вне параболы $b^2 = 4\Delta$ (т.е. $b^2 > 4\Delta$), соответствуют точкам покоя типа узла.

Область $\Delta < 0$ пл-ти $bO\Delta$ --- точки типа седла.

$\mathcal{D} = b^2 - 4\Delta > 0 \implies$ корни вещественные по т. Виета.

Рассмотрим систему линейных однородных дифференциальных уравнений с постоянными коэффициентами:
\begin{equation} \label{eq:231}
	\dv{x_i}{t} = \sum_{j=1}^{n} a_{ij} x_j, \quad i=1,2,\dots,n \quad (n \ge 2) \quad (231)
\end{equation}

Для нее имеют место аналогичные типы расположения интегральных кривых около начала координат (обобщенное седло, обобщенный узел).

\vspace{0.5cm}
\begin{center}
	\begin{tikzpicture}[scale=0.9, decoration={markings, mark=at position 1 with {\arrow{stealth}}}]
		% n=1
		\begin{scope}[xshift=-4cm, yshift=2cm]
			\node[left] at (-1.5, 1.5) {$n=1$};
			\draw[-stealth, thick] (-1, 0) -- (3, 0) node[right] {$\mathbb{R}$};
			\filldraw (0.5, 0) circle (1.5pt) node[below] {$0$};
		\end{scope}
		
		% n=3
		\begin{scope}[xshift=2cm, yshift=2cm]
			\node[left] at (-1.5, 1.5) {$n=3$};
			\draw[-stealth, thick] (0, 0) -- (0, 2) node[right] {$x_3$};
			\draw[-stealth, thick] (0, 0) -- (2, 0) node[right] {$x_2$};
			\draw[-stealth, thick] (0, 0) -- (-1.5, -1.5) node[below] {$x_1$};
		\end{scope}
		
		% n=4
		\begin{scope}[xshift=8cm, yshift=2cm]
			\node[left] at (-1.5, 1.5) {$n=4$};
			\node[above, align=center] at (1, 2) {пр-во Минковского};
			\draw[-stealth, thick] (0, -1) -- (0, 1.5) node[right] {$t$};
			\draw[-stealth, thick] (0, -1) -- (2.5, -1) node[below] {$x_1, x_2, x_3$};
			\node[below left] at (0, -1) {$0$};
		\end{scope}
		
		% n=2
		\begin{scope}[xshift=-4cm, yshift=-2.5cm]
			\node[left] at (-1.5, 1.5) {$n=2$};
			\draw[-stealth, thick] (0, -0.5) -- (0, 2) node[left] {$y$};
			\draw[-stealth, thick] (-0.5, 0) -- (2.5, 0) node[right] {$x$};
			\node[below left] at (0, 0) {$0$};
			\node at (1.5, 1) {$\mathbb{R}^2$};
		\end{scope}
	\end{tikzpicture}
\end{center}
\vspace{0.5cm}

\begin{theorem}{Об устойчивости точки покоя системы линейных ДУ}{linear_stability}
	Если все корни характеристического уравнения для $n$-мерной системы имеют отрицательную вещественную часть, то точка покоя системы \eqref{eq:231} $x_i = 0$, $i=1,2,\dots,n$, асимптотически устойчива. Если хотя бы один корень характеристического уравнения имеет положительную вещественную часть, то точка покоя неустойчива.
\end{theorem}

\begin{example}
	Проверим устойчивость точки покоя $(0,0,0)$ системы:
	\begin{equation*}
		\begin{cases}
			\dv{x}{t} = -x + z \\
			\dv{y}{t} = -2y - z \\
			\dv{z}{t} = y - z
		\end{cases}
	\end{equation*}
	Составляем характеристическое уравнение:
	\begin{equation*}
		\begin{vmatrix} 
			-1-\lambda & 0 & 1 \\ 
			0 & -2-\lambda & -1 \\ 
			0 & 1 & -1-\lambda 
		\end{vmatrix} = 0
	\end{equation*}
	\begin{equation*}
		\Updownarrow
	\end{equation*}
	\begin{equation*}
		-(1+\lambda)(2+\lambda+2\lambda+\lambda^2+1) = 0
	\end{equation*}
	\begin{equation*}
		\Updownarrow
	\end{equation*}
	\begin{equation*}
		(1+\lambda)(\lambda^2+3\lambda+3) = 0
	\end{equation*}
	\begin{equation*}
		\lambda_1 = -1, \quad \lambda_{2,3} = \frac{-3 \pm \sqrt{-3}}{2} = -\frac{3}{2} \pm i\frac{\sqrt{3}}{2}
	\end{equation*}
	\begin{center}
		все корни имеют отрицательные вещественные части
	\end{center}
	\begin{equation*}
		\Downarrow
	\end{equation*}
	\begin{center}
		точка покоя $(0,0,0)$ --- асимптотически устойчива
	\end{center}
\end{example}

\section{Метод функций Ляпунова}

Метод функций Ляпунова состоит в непосредственном исследовании устойчивости положения равновесия системы $\dv{x_i}{t} = f_i(x_1, x_2, \dots, x_n)$, $i=1,2,\dots,n$ при помощи подходящим образом подобранной функции $V(t, x_1, x_2, \dots, x_n)$ --- функции Ляпунова, причем это делается без предварительного нахождения решений системы.

\begin{itemize}
	\item \textbf{Автономные системы:}
\end{itemize}
\begin{equation} \label{eq:232}
	\dv{x_i}{t} = f_i(x_1, x_2, \dots, x_n), \quad i=1,2,\dots,n \quad (232)
\end{equation}
для которых $x_i \equiv 0$, $i=1,2,\dots,n$ есть точка покоя.

\begin{definition}{Знакоопределенная функция}{def_positive_definite}
	Функция $V(x_1, x_2, \dots, x_n)$, определенная в некоторой окрестности начала координат, называется \textbf{знакоопределенной} (определенно-положительной / определенно-отрицательной), если она в области 
	\begin{equation} \label{eq:233}
		|x_i| \le h, \quad i=1,2,\dots,n \quad (233)
	\end{equation}
	где $h$ --- достаточно малое положительное число, может принимать значения только одного определенного знака и обращается в $0$ лишь при $x_1 = x_2 = \dots = x_n = 0$.
\end{definition}

Например, функции $V = x_1^2 + x_2^2 + x_3^2$ и $V = x_1^2 + 2x_1x_2 + 2x_2^2 + x_3^2$ будут определенно-положительными, причем величина $h > 0$ может быть взята сколь угодно большой.

\begin{definition}{Знакопостоянная функция}{def_positive_semidefinite}
	Функция $V(x_1, x_2, \dots, x_n)$ называют \textbf{знакопостоянной} (положительной / отрицательной), если она в области $|x_i| \le h$, $i=1,2,\dots,n$ может принимать значения только одного определенного знака, но может обращаться в $0$ и при $x_1^2 + x_2^2 + \dots + x_n^2 \neq 0$.
\end{definition}

Например, функция $V = x_1^2 + x_2^2 + 2x_1x_2 + x_3^2 = (x_1+x_2)^2 + x_3^2$ будет знакопостоянной (положительной), но не будет знакоопределенной, так как она обращается в $0$ и при $x_1^2 + x_2^2 + x_3^2 \neq 0$, а именно:
$\begin{cases} x_3 = 0 \\ x_1 = -x_2 \end{cases}$.

\begin{definition}{Полная производная функции по времени}{def_total_derivative}
	Пусть $V(x_1, x_2, \dots, x_n)$ есть дифференцируемая функция своих аргументов. Пусть $x_1, x_2, \dots, x_n$ --- являются некоторыми функциями времени, удовлетворяющими системе дифференциальных уравнений \eqref{eq:232}. Тогда
	\begin{equation} \label{eq:234}
		\dv{V}{t} = \sum_{i=1}^{n} \pdv{V}{x_i} \cdot \dv{x_i}{t} = \sum_{i=1}^{n} \pdv{V}{x_i} \cdot f_i(x_1, \dots, x_n) \quad (234)
	\end{equation}
	Величина $\dv{V}{t}$, определяемая формулой \eqref{eq:234}, называется \textbf{полной производной функции $V$ по времени}, составленной в силу системы уравнений \eqref{eq:232}.
\end{definition}

\begin{theorem}{Теорема Ляпунова об устойчивости}{lyap_stability}
	Если для системы дифференциальных уравнений \eqref{eq:232} существует знакоопределенная функция $V(x_1, x_2, \dots, x_n)$ (функция Ляпунова), полная производная $\dv{V}{t}$ которой по времени, составленная в силу системы \eqref{eq:232}, есть функция знакопостоянная, знака, противоположного с $V$, или тождественно равная $0$, то точка покоя $x_i = 0$, $i=1,2,\dots,n$, системы \eqref{eq:232} устойчива.
\end{theorem}

\begin{theorem}{Теорема Ляпунова об асимптотической устойчивости}{lyap_asymp_stability}
	Если для системы дифференциальных уравнений \eqref{eq:232} существует знакоопределенная функция $V(x_1, x_2, \dots, x_n)$ (функция Ляпунова), полная производная $\dv{V}{t}$ которой по времени, составленная в силу системы \eqref{eq:232}, есть функция знакоопределенная знака, противоположного с $V$, то точка покоя $x_i = 0$, $i=1,2,\dots,n$, системы \eqref{eq:232} асимптотически устойчива.
\end{theorem}

\begin{example}
	Исследовать на устойчивость точку покоя системы:
	\begin{equation} \label{eq:235}
		\begin{cases}
			\dv{x}{t} = y \\
			\dv{y}{t} = -x
		\end{cases} \quad (235)
	\end{equation}
	Выберем в качестве функции $V(x,y)$ функцию $V = x^2 + y^2$.
	Эта функция определенно-положительная.
	Производная функции $V$ в силу системы \eqref{eq:235} равна:
	\begin{equation*}
		\dv{V}{t} = \underbrace{\pdv{V}{x}}_{2x} \cdot \underbrace{\dv{x}{t}}_{y} + \underbrace{\pdv{V}{y}}_{2y} \cdot \underbrace{\dv{y}{t}}_{-x} = 2xy - 2xy \equiv 0.
	\end{equation*}
	Однако асимптотической устойчивости нет: траектории системы \eqref{eq:235} --- окружности, и они не стремятся к точке $(0,0)$ при $t \to +\infty$.
	$\Downarrow$
	по Теореме Ляпунова об устойчивости, точка покоя $(0,0)$ системы \eqref{eq:235} устойчива.
\end{example}

\begin{example}
	Исследовать на устойчивость точку покоя системы:
	\begin{equation} \label{eq:236}
		\begin{cases}
			\dv{x}{t} = y - x^3 \\
			\dv{y}{t} = -x - 3y^3
		\end{cases} \quad (236)
	\end{equation}
	Возьмем $V = x^2 + y^2$, тогда:
	\begin{align*}
		\dv{V}{t} &= \underbrace{2x}_{\pdv{V}{x}} \underbrace{(y - x^3)}_{\dv{x}{t}} + \underbrace{2y}_{\pdv{V}{y}} \underbrace{(-x - 3y^3)}_{\dv{y}{t}} \\
		&= -2(x^4 + 3y^4) < 0 \implies \text{определенно-отриц. функция}
	\end{align*}
	$\Downarrow$
	по Теореме Ляпунова об асимптотической устойчивости, точка покоя $(0,0)$ системы \eqref{eq:236} устойчива асимптотически.
\end{example}

\begin{nb}
	Общего метода построения функции Ляпунова \textbf{НЕТ}. В простейших случаях функцию Ляпунова можно искать в виде: $V(x,y) = ax^2 + by^2$, $V(x,y) = ax^4 + bx^4$, $V(x,y) = ax^4 + by^4$ ($a>0, b>0$) и так далее.
\end{nb}

\begin{example}
	С помощью функции Ляпунова исследовать на устойчивость тривиальное решение $x=0, y=0$ системы:
	\begin{equation*}
		\begin{cases}
			\dv{x}{t} = -x - 2y + x^2y^2 \\
			\dv{y}{t} = x - \frac{y}{2} - \frac{x^3y}{2}
		\end{cases}
	\end{equation*}
	Будем искать функцию Ляпунова в виде: $V = ax^2 + by^2$, где $a>0, b>0$.
	\begin{align*}
		\dv{V}{t} = \pdv{V}{x}\dv{x}{t} + \pdv{V}{y}\dv{y}{t} &= 2ax(-x - 2y + x^2y^2) + 2by\qty(x - \frac{y}{2} - \frac{x^3y}{2}) \\
		&= -(2ax^2 + by^2) + (2xy - x^3y^2)(b - 2a).
	\end{align*}
	Положим $b = 2a$, тогда
	\begin{align*}
		\dv{V}{t} &= -2a(x^2 + y^2) \le 0 \\
		\implies &\text{при любом } a>0 \text{ и } b=2a \text{ функция } V = ax^2 + 2ay^2 \\
		&\text{будет опред.-положительной},
	\end{align*}
	а $\dv{V}{t}$ --- определенно-отрицательной.
	$\implies x=0, y=0$ устойчиво асимптотически.
\end{example}

\begin{theorem}{Теорема Ляпунова о неустойчивости}{lyap_instability}
	Пусть для системы дифференциальных уравнений \eqref{eq:232} существует дифференцируемая в окрестности начала координат функция $V(x_1, x_2, \dots, x_n)$ такая, что $V(0,0,\dots,0) = 0$. Если ее полная производная $\dv{V}{t}$, составленная в силу системы \eqref{eq:232}, есть определенно-положительная функция и сколь угодно близко от начала координат имеются точки, в которых функция $V(x_1, x_2, \dots, x_n) > 0$, то точка покоя $x_i = 0$ неустойчива.
\end{theorem}

\section{Устойчивость по первому приближению}

\begin{itemize}
	\item Рассмотрим систему дифференциальных уравнений:
\end{itemize}
\begin{equation} \label{eq:237}
	\dv{x_i}{t} = f_i(x_1, x_2, \dots, x_n), \quad i=1,2,\dots,n \quad (237)
\end{equation}
Пусть $x_i \equiv 0$, $i=1,2,\dots,n$ есть точка покоя системы \eqref{eq:237}, то есть $f_i(0, 0, \dots, 0) = 0$, $i=1,2,\dots,n$.
Предположим, что $f_i(x)$ дифференцируемы в начале координат достаточное число раз.

Разложим функции $f_i$ по формуле Тейлора по $x$ в окрестности начала координат:
\begin{equation*}
	f_i(x_1, x_2, \dots, x_n) = \sum_{j=1}^{n} a_{ij} x_j + R_i(x_1, x_2, \dots, x_n), \quad \text{где } a_{ij} = \pdv{f_i}{x_j}(0, 0, \dots, 0),
\end{equation*}
а $R_i$ --- члены второго порядка малости относительно начала координат.

Тогда система \eqref{eq:237} запишется в следующем виде:
\begin{equation} \label{eq:238}
	\dv{x_i}{t} = \sum_{j=1}^{n} a_{ij} x_j + R_i(x_1, x_2, \dots, x_n), \quad i=1,2,\dots,n \quad (238)
\end{equation}

Вместо системы \eqref{eq:238} рассмотрим систему:
\begin{equation} \label{eq:239}
	\dv{x_i}{t} = \sum_{j=1}^{n} a_{ij} x_j, \quad i=1,2,\dots,n, \quad a_{ij} = \text{const} \quad (239)
\end{equation}
называемую \textbf{системой уравнений первого приближения} для системы \eqref{eq:237}.

\begin{theorem}{Утверждение}{stability_first_approx}
	1) Если все корни характеристического уравнения
	\begin{equation} \label{eq:240}
		\begin{vmatrix}
			a_{11} - \lambda & a_{12} & \dots & a_{1n} \\
			a_{21} & a_{22} - \lambda & \dots & a_{2n} \\
			\dots & \dots & \dots & \dots \\
			a_{n1} & a_{n2} & \dots & a_{nn} - \lambda
		\end{vmatrix} = 0 \quad (240)
	\end{equation}
	имеют отрицательные вещественные части, то нулевое решение $x_i = 0$, $i=1,2,\dots,n$ системы \eqref{eq:239} и системы \eqref{eq:238} асимптотически устойчиво.
	
	2) Если все корни характеристического уравнения \eqref{eq:240} имеют положительную вещественную часть, то нулевое решение системы \eqref{eq:239} и системы \eqref{eq:238} неустойчиво.
\end{theorem}

Говорят, что в случаях 1 и 2 возможно исследование решения на устойчивость по первому приближению. Случаи, когда вещественные части всех корней характеристического уравнения \eqref{eq:240} неположительны (то есть отрицательные или нулевые), являются критическими.
Если вещественная часть хотя бы одного корня равна $0$, то исследование решения на устойчивость по первому приближению, вообще говоря, невозможно (так как начинают влиять нелинейные члены $R_i$).

\begin{example}
	Исследовать на устойчивость по первому приближению точку покоя $(0,0)$ системы:
	\begin{equation} \label{eq:241}
		\begin{cases}
			\dot{x} = 2x + y - 5y^2 \\
			\dot{y} = 3x + y + \frac{x^3}{2}
		\end{cases} \quad (241)
	\end{equation}
	Система первого приближения:
	\begin{equation} \label{eq:242}
		\begin{cases}
			\dot{x} = 2x + y \\
			\dot{y} = 3x + y
		\end{cases} \quad (242)
	\end{equation}
	Составим характеристическое уравнение:
	\begin{equation} \label{eq:243}
		\det \begin{pmatrix} 2-\lambda & 1 \\ 3 & 1-\lambda \end{pmatrix} = 0 \iff (2-\lambda)(1-\lambda) - 3 = 0 \iff \lambda^2 - 3\lambda - 1 = 0 \quad (243)
	\end{equation}
	\begin{equation*}
		D = 9 + 4 = 13 \implies \lambda_{1,2} = \frac{3 \pm \sqrt{13}}{2} \text{ --- вещественные корни}
	\end{equation*}
	$\lambda_1 > 0 \implies$ нулевое решение системы \eqref{eq:241} неустойчиво (по п. 2 утверждения).
\end{example}

\begin{example}
	Исследовать на устойчивость по первому приближению точку покоя $x=0, y=0$ систем:
	\begin{equation} \label{eq:244}
		\begin{cases}
			\dv{x}{t} = y - x^3 \\
			\dv{y}{t} = -x - y^3
		\end{cases} \quad (244)
	\end{equation}
	и
	\begin{equation} \label{eq:245}
		\begin{cases}
			\dv{x}{t} = y + x^3 \\
			\dv{y}{t} = -x + y^3
		\end{cases} \quad (245)
	\end{equation}
	Проверим устойчивость решения методом Ляпунова. Для системы \eqref{eq:244} возьмем функцию Ляпунова: $V = x^2 + y^2$.
	\begin{equation*}
		\dv{V}{t} = \underbrace{\pdv{V}{x}}_{2x} \underbrace{\dv{x}{t}}_{y-x^3} + \underbrace{\pdv{V}{y}}_{2y} \underbrace{\dv{y}{t}}_{-x-y^3} = 2x(y-x^3) + 2y(-x-y^3) = -2(x^4+y^4) \le 0
	\end{equation*}
	$\implies$ по Теореме Ляпунова точка покоя $x=0, y=0$ системы \eqref{eq:244} асимптотически устойчива.
	
	Для системы \eqref{eq:245} возьмем $V = x^2 + y^2$:
	\begin{equation*}
		\dv{V}{t} = 2x(y+x^3) + 2y(-x+y^3) = 2(x^4+y^4) > 0
	\end{equation*}
	$\implies$ по Теореме Ляпунова точка покоя $x=0, y=0$ системы \eqref{eq:245} неустойчива.
	
	Системы \eqref{eq:244} и \eqref{eq:245} имеют одну и ту же систему первого приближения:
	\begin{equation} \label{eq:246}
		\begin{cases}
			\dv{x}{t} = y \\
			\dv{y}{t} = -x
		\end{cases} \quad (246)
	\end{equation}
	Характеристическое уравнение: $\begin{vmatrix} -\lambda & 1 \\ -1 & -\lambda \end{vmatrix} = 0 \iff \lambda^2 + 1 = 0 \iff \lambda^2 = -1 \iff \lambda = \pm i$ --- чисто мнимые корни, то есть вещественные части корней характеристического уравнения равны $0$.
	
	Для системы первого приближения \eqref{eq:246} начало координат является центром. Системы \eqref{eq:244} и \eqref{eq:245} получаются малым возмущением правых частей системы \eqref{eq:246} в окрестности начала координат.
	Однако эти малые возмущения приводят к тому, что замкнутые траектории превращаются в спирали, в случае \eqref{eq:244} приближающиеся к началу координат и образующие в точке $(0,0)$ устойчивый фокус, а в случае \eqref{eq:245} --- удаляющиеся от начала координат и образующие в точке $(0,0)$ неустойчивый фокус.
	
	\textbf{Вывод:} В критическом случае нелинейные члены могут влиять на устойчивость точки покоя.
\end{example}

\begin{example}
	Замкнутый контур с нелинейными элементами.
	
	\begin{center}
		\begin{circuitikz}[european, scale=1.2]
			\draw (0,0) -- (0,2) 
			to[L, l_=$L$] (3,2) 
			to[R=$R$] (6,2) 
			to[C=$C$] (6,0) 
			-- (0,0);
		\end{circuitikz}
	\end{center}
	
	Уравнение контура:
	\begin{equation} \label{eq:247}
		L \dv[2]{x}{t} + R \dv{x}{t} + \frac{1}{C} x + g\qty(x, \dv{x}{t}) = 0 \quad (247)
	\end{equation}
	Здесь $x$ --- заряд конденсатора $\implies \dv{x}{t}$ --- ток в цепи, $R$ --- сопротивление, $L$ --- индуктивность, $C$ --- ёмкость. $g\qty(x, \dv{x}{t})$ --- нелинейные члены, имеющие степень не ниже второй, $g(0,0) = 0$.
	
	Перепишем \eqref{eq:247} в виде:
	\begin{equation} \label{eq:248}
		\begin{cases}
			\dv{x}{t} = y \\
			\dv{y}{t} = -\frac{1}{LC} x - \frac{R}{L} y - \frac{1}{L} g(x,y)
		\end{cases} \quad (248)
	\end{equation}
	
	Запишем систему первого приближения:
	\begin{equation} \label{eq:249}
		\begin{cases}
			\dv{x}{t} = y \\
			\dv{y}{t} = -\frac{1}{LC} x - \frac{R}{L} y
		\end{cases} \quad (249)
	\end{equation}
	Характеристическое уравнение для \eqref{eq:249}:
	\begin{equation} \label{eq:250}
		\det \begin{pmatrix} -\lambda & 1 \\ -\frac{1}{LC} & -\frac{R}{L}-\lambda \end{pmatrix} = 0 \iff \lambda^2 + \frac{R}{L}\lambda + \frac{1}{LC} = 0 \iff \lambda = \frac{-\frac{R}{L} \pm \sqrt{\frac{R^2}{L^2} - \frac{4}{LC}}}{2} \quad (250)
	\end{equation}
	$\iff \lambda = -\frac{R}{2L} \pm \sqrt{\frac{R^2}{4L^2} - \frac{1}{LC}}$. 
	Если $D < 0$, то есть $\frac{R^2}{4L^2} < \frac{1}{LC} \iff R^2 < \frac{4L}{C}$, то уравнение \eqref{eq:250} имеет комплексные корни с отрицательной вещественной частью $\qty(-\frac{R}{2L}) \implies$ точка $(0,0)$ системы \eqref{eq:249} и \eqref{eq:248} асимптотически устойчива. Если $D > 0$, то есть $R^2 > \frac{4L}{C}$, то точка $(0,0)$ асимптотически устойчива, так как $\lambda < 0$.
\end{example}

\section{Устойчивость решений при изменении правых частей}

\begin{itemize}
	\item Дифференциальные уравнения:
\end{itemize}
\begin{align}
	y' &= f(x,y) \quad (251) \label{eq:251} \\
	y' &= f(x,y) + \theta(x,y) \quad (252) \label{eq:252}
\end{align}
где функции $f(x,y)$ и $\theta(x,y)$ непрерывны в замкнутой области $\overline{G}$ плоскости $xOy$ и функция $f(x,y)$ имеет в этой области непрерывную частную производную $\pdv{f}{y}$.

\begin{theorem}{Об оценке разности решений при малом изменении правой части уравнения}{diff_estimation}
	Пусть в области $\overline{G}$ выполняется неравенство $|\theta(x,y)| \le \varepsilon$.
	Если $y = \varphi(x)$ и $y = \psi(x)$ есть решения уравнений \eqref{eq:251} и \eqref{eq:252} соответственно, удовлетворяющие одному и тому же начальному условию $\left. \varphi \right|_{x=x_0} = \left. \psi \right|_{x=x_0} = y_0$, то:
	\begin{equation} \label{eq:253}
		|\varphi(x) - \psi(x)| \le \frac{\varepsilon}{M} \qty( \exp(M |x - x_0|) - 1 ), \quad \text{где } M = \max_{(x,y) \in \overline{G}} \abs{\pdv{f}{y}} \quad (253)
	\end{equation}
\end{theorem}
